\section*{Chapitre \ref{ch:b1:stress-strain-rheology} --- Contrainte, déformation et rhéologie}

\begin{solution}{exo:b1:stress-strain-rheology:1}
(a) $\sigma_v = 2700\times9.81\times10\,000 = \qty{2.65e8}{Pa} = \qty{264.9}{MPa}$.
(b) $P = 1000\times9.81\times10\,000 = \qty{98.1}{MPa}$. (c) $\sigma_v' =
264.9 - 98.1 = \qty{166.8}{MPa}$, et $\lambda = 98.1/264.9 = 0.370$, ce qui vaut
$1000/2700$.
\end{solution}

\begin{solution}{exo:b1:stress-strain-rheology:2}
(a) $e = (13-10)/10 = 0.30$ et $S = 1.30$. (b) $e = (7-10)/10 = -0.30$, un
raccourcissement de \qty{30}{\percent}, et $S = 0.70$. (c) L'angle s'est fermé de
$\psi = \ang{30}$, et $\gamma = \tan\ang{30} = 0.577$.
\end{solution}

\begin{solution}{exo:b1:stress-strain-rheology:3}
(a) Centre \qty{70}{MPa}, rayon \qty{30}{MPa}. (b) La plus grande contrainte de cisaillement
est le rayon, \qty{30}{MPa}, sur les plans à \ang{45} de $\sigma_1$ et de
$\sigma_3$, qui portent la \omterm{def:b1:stress-strain-rheology:stress}{contrainte normale} du centre, \qty{70}{MPa}. (c)
$2\theta = \ang{60}$~: $\sigma_n = 70 + 30\cos\ang{60} = \qty{85}{MPa}$ et
$\tau = 30\sin\ang{60} = \qty{26.0}{MPa}$.
\end{solution}

\begin{solution}{exo:b1:stress-strain-rheology:4}
(a) $\ang{45} - \ang{30}/2 = \ang{30}$ avec $\sigma_1$, de part et d'autre d'elle,
donc les deux plans font \ang{60} entre eux. (b) Avec $c' = 20 + r$, la
condition du \cref{thm:b1:stress-strain-rheology:failure-angle} s'écrit
$r = C\cos\phi + (20 + r)\sin\phi$, donc
\[
  r = \frac{C\cos\phi + 20\sin\phi}{1 - \sin\phi} = \frac{8.66 + 10}{0.5}
    = \qty{37.3}{MPa},
\]
et $\sigma_1' = 20 + 2r = \qty{94.6}{MPa}$.
\end{solution}

\begin{solution}{exo:b1:stress-strain-rheology:5}
(a) $c = 40$ et $R = \sqrt{20^2 + 15^2} = \qty{25}{MPa}$~: $\sigma_1 =
\qty{65}{MPa}$ et $\sigma_3 = \qty{15}{MPa}$~; $\tan2\theta_p = 30/40 = 0.75$,
$2\theta_p = \ang{36.87}$, $\theta_p = \ang{18.43}$. (b) À cet angle
$\cos^2\theta = 0.9$, $\sin^2\theta = 0.1$ et $\sin\theta\cos\theta = 0.3$~:
$\sigma_n = 60\times0.9 + 2\times15\times0.3 + 20\times0.1 = 54 + 9 + 2 =
\qty{65}{MPa}$, et $\tau = (20-60)\times0.3 + 15\times(0.9-0.1) = -12 + 12 =
0$. (c) $R = \qty{25}{MPa}$, sur les plans à \ang{45} des directions
principales.
\end{solution}

\begin{solution}{exo:b1:stress-strain-rheology:6}
(a) $S_{1,2} = \sqrt{1 + 1} \pm 1$~: $S_1 = 2.414$ et $S_2 = 0.414$~;
$\tan2\theta' = 2/2 = 1$, $\theta' = \ang{22.5}$. (b) $S_1S_2 =
(\sqrt2 + 1)(\sqrt2 - 1) = 1$, qui est $\det F$~: pas de changement d'aire. (c) La
ligne selon $(1, 1)$ va en $F(1, 1) = (1 + 2, 1) = (3, 1)$, à
$\arctan(1/3) = \ang{18.4}$ de la direction de cisaillement, avec l'\omterm{def:b1:stress-strain-rheology:strain}{étirement}
$\sqrt{10}/\sqrt2 = 2.24$. Elle a dépassé le grand axe de l'ellipse,
à \ang{22.5}, vers le plan de cisaillement~; son \omterm{def:b1:stress-strain-rheology:strain}{étirement} est un peu inférieur à
$S_1$, comme il se doit, puisqu'aucune ligne ne s'étire plus que $S_1$.
\end{solution}

\begin{solution}{exo:b1:stress-strain-rheology:7}
(a) $S_1/S_2 = 4$ et $S_1S_2 = 1$ donnent $S_1 = 2$ et $S_2 = 0.5$~; d'après la
\cref{prop:b1:stress-strain-rheology:simple-shear}, $\gamma = S_1 - S_2 =
1.5$. (b) $\tan2\theta' = 2/1.5 = 1.333$, $2\theta' = \ang{53.1}$~: les grands
axes devraient se coucher à \ang{26.6} des parois. (c) $1.5\times200 = \qty{300}{m}$.
\end{solution}

\begin{solution}{exo:b1:stress-strain-rheology:8}
$\sigma_v = 2700\times9.81\times3000 = \qty{79.46}{MPa}$. (a) $P = \qty{29.43}{MPa}$,
$\sigma_n' = 70 - 29.43 = \qty{40.57}{MPa}$, et Byerlee permet
$0.85\times40.57 = \qty{34.5}{MPa}$, plus que les \qty{30}{MPa} appliqués~: la
faille tient. (b) $P = 0.6\times79.46 = \qty{47.68}{MPa}$, $\sigma_n' =
\qty{22.32}{MPa}$, et le frottement ne donne que \qty{19.0}{MPa}~: la faille
glisse. (c) $30 = 0.85\,(70 - 79.46\,\lambda)$ donne $\lambda = (70 -
35.29)/79.46 = 0.437$.
\end{solution}

\begin{solution}{exo:b1:stress-strain-rheology:9}
(a) La contrainte varie comme $\dot\varepsilon^{1/4}$~: $10^{1/4} = 1.78$. (b) Le
rapport des résistances est
\[
  \exp\Bigl[\frac{E}{nR}\Bigl(\frac{1}{673.15} - \frac{1}{773.15}\Bigr)\Bigr]
  = \exp\bigl(6705\times\num{1.921e-4}\bigr) = \exp(1.288) = 3.63 .
\]
(c) Un changement de vitesse d'un facteur dix change la résistance au fluage d'un facteur 1.78,
ce que la température fait en seulement $100\times\ln1.78/\ln3.63 \approx
\qty{45}{\degreeCelsius}$~; sur un \omterm{def:b1:heat-budget:geotherm}{géotherme} d'environ \qty{20}{\degreeCelsius}
par kilomètre, cela déplace la transition d'environ \qty{2}{km}. Un changement de
\omterm{def:b1:heat-budget:geotherm}{géotherme} déplace les isothermes, et avec elles la transition, de bien davantage.
\end{solution}

\begin{solution}{exo:b1:stress-strain-rheology:10}
(a) Avec $C = 0$, une faille glisse quand $r = \mu c'/\sqrt{1+\mu^2}$
(\cref{thm:b1:stress-strain-rheology:failure-angle}). Ici $c' = S =
\rho g z(1-\lambda)$, donc $\sigma_1 - \sigma_3 = 2r = 2\mu S/\sqrt{1+\mu^2}$.
(b) $S$ croît de $2700\times9.81\times1000\times(1 - 1000/2700) =
\qty{16.68}{MPa}$ par kilomètre, et $2\times0.85/1.3124 = 1.295$~: la pente
est de \qty{21.6}{MPa} par kilomètre, entre \num{13.1} (failles normales) et
\qty{61.3}{MPa} par kilomètre (chevauchements). (c) $2\mu/(k+\mu) = 0.786$,
$2\mu/k = 1.295$ et $2\mu/(k-\mu) = 3.68$, avec $k = 1.3124$~: la
\omterm{def:b1:stress-strain-rheology:principal}{contrainte différentielle} au glissement vaut 0.79, 1.30 et 3.68 fois la \omterm{def:b1:stress-strain-rheology:lithostatic}{contrainte
effective} verticale.
\end{solution}

\begin{solution}{exo:b1:stress-strain-rheology:11}
(a) $t = 60\times\num{e6}\times\num{3.156e7} = \qty{1.89e15}{s}$ et
$2\sqrt{\kappa t} = 2\sqrt{\num{1.89e9}} = \qty{87.0}{km}$. (b)
$\operatorname{erf}x = 600/1350 = 0.444$ pour $x = 0.417$~: $z = 0.417\times87.0
= \qty{36.3}{km}$~; $\operatorname{erf}x = 800/1350 = 0.593$ pour $x = 0.586$~:
$z = \qty{51.0}{km}$. (c) Les séismes devraient atteindre \qtyrange{36}{51}{km}~;
l'enveloppe place la base de la couche cassante à \qty{38.2}{km}, dans
cet intervalle et près de son sommet, à une température, \qty{641}{\degreeCelsius}, dans
l'intervalle observé.
\end{solution}

\begin{solution}{exo:b1:stress-strain-rheology:12}
(a) $r = 45$ et $c' = \qty{65}{MPa}$~; la distance du centre à la droite
$\tau = 0.85\sigma_n'$ est $0.85\times65/1.3124 = \qty{42.1}{MPa}$, inférieure au
rayon~: le cercle coupe la droite, et les plans sur l'arc au-dessus d'elle
glissent. (b) Avec $2\theta$ l'angle depuis l'extrémité $\sigma_1$, le glissement demande
$45\sin2\theta \ge 0.85\,(65 + 45\cos2\theta)$, c'est-à-dire
$45\sin2\theta - 38.25\cos2\theta \ge 55.25$, ou
$59.06\sin(2\theta - \ang{40.36}) \ge 55.25$, où $\ang{40.36} = \phi$.
Donc $\sin(2\theta - \phi) \ge 0.9355$, $2\theta - \phi$ entre \ang{69.31}
et \ang{110.69}, $2\theta$ entre \ang{109.67} et \ang{151.06}~: les normales
se trouvent à \ang{54.8} à \ang{75.5} de $\sigma_1$, et les fractures qui glissent
à \ang{14.5} à \ang{35.2} de $\sigma_1$. Le milieu de l'intervalle,
\ang{24.8}, est $\ang{45} - \phi/2$, le plan du
\cref{thm:b1:stress-strain-rheology:failure-angle}. (c) Seul ce plan glisse
quand le cercle touche juste la droite~: $r = \mu(20 + r)/k$, $r = 20\mu/(k -
\mu) = 17/0.4624 = \qty{36.8}{MPa}$, et $\sigma_1' = 20 + 2r =
\qty{93.5}{MPa}$.
\end{solution}

\begin{solution}{pb:b1:stress-strain-rheology:1}
\textbf{1.} $\sigma_v = 2700\times9.81\times5000 = \qty{132.44}{MPa}$.

\textbf{2.} $P = 1000\times9.81\times5000 = \qty{49.05}{MPa}$, et
$\lambda = 49.05/132.44 = 0.370$.

\textbf{3.} $\sigma_1 = \sigma_v = \qty{132.44}{MPa}$, verticale, et
$\sigma_3 = \qty{75.00}{MPa}$, horizontale est--ouest~: la \omterm{def:b1:layered-earth:layers}{croûte} est étirée
est--ouest sous son propre poids, l'état qui entraîne des failles
normales, comme celle-ci.

\textbf{4.} $\sigma_1 - \sigma_3 = \qty{57.44}{MPa}$~; centre
$(132.44 + 75.00)/2 = \qty{103.72}{MPa}$~; rayon \qty{28.72}{MPa}.

\textbf{5.} $\sigma_1' = 132.44 - 49.05 = \qty{83.39}{MPa}$, $\sigma_3' =
75.00 - 49.05 = \qty{25.95}{MPa}$, centre $103.72 - 49.05 = \qty{54.67}{MPa}$.

\textbf{6.} $\sigma_{xx} = \qty{75.00}{MPa}$, $\sigma_{yy} = \qty{132.44}{MPa}$
et $\sigma_{xy} = 0$~: les axes sont les directions principales, l'une d'elles
verticale, donc les plans perpendiculaires à elles ne portent aucune contrainte de cisaillement.

\textbf{7.} Selon son \omterm{def:b1:geological-maps:attitude}{pendage} la faille suit le vecteur unitaire
$(\cos\ang{60}, -\sin\ang{60})$, vers l'est et vers le bas~; sa normale vers le haut est
$\vec n = (\sin\ang{60}, \cos\ang{60}) = (0.866, 0.500)$, à \ang{30} de
$x$. Son angle avec $\sigma_1$, selon $y$, a pour cosinus
0.500~: $\theta = \ang{60}$.

\textbf{8.} Avec \ang{30} de $x$~: $\sigma_n = 75.00\times0.75 + 132.44\times0.25
= 56.25 + 33.11 = \qty{89.36}{MPa}$, et $\tau = (132.44 - 75.00)\times
\sin\ang{30}\cos\ang{30} = 57.44\times0.4330 = \qty{24.87}{MPa}$.

\textbf{9.} À $2\theta = \ang{120}$ de l'extrémité $\sigma_1$~:
$\sigma_n = 103.72 + 28.72\cos\ang{120} = 103.72 - 14.36 = \qty{89.36}{MPa}$
et $\tau = 28.72\sin\ang{120} = \qty{24.87}{MPa}$.

\textbf{10.} $\sigma_n' = 89.36 - 49.05 = \qty{40.31}{MPa}$.

\textbf{11.} $24.87/40.31 = 0.617$, inférieur à 0.85~: la faille tient.

\textbf{12.} $0.85\times40.31 = \qty{34.26}{MPa}$~; la contrainte de cisaillement est
en deçà de \qty{9.39}{MPa}.

\textbf{13.} Le cercle a son centre à \qty{54.67}{MPa} et pour rayon
\qty{28.72}{MPa}, de \num{25.95} à \qty{83.39}{MPa}~; la droite de Byerlee est
$\tau = 0.85\,\sigma_n'$ (toutes les \omterm{def:b1:stress-strain-rheology:stress}{contraintes normales} ici sont inférieures à
\qty{200}{MPa})~; la faille est le point $(40.31, 24.87)$, à \ang{120} de
l'extrémité $\sigma_1$.

\textbf{14.} $0.85\times54.67/\sqrt{1 + 0.85^2} = 46.47/1.3124 =
\qty{35.41}{MPa}$, plus que le rayon, \qty{28.72}{MPa}~: tout le cercle
est sous la droite, et aucun plan de la roche ne peut glisser.

\textbf{15.} $\phi = \arctan0.85 = \ang{40.36}$~: le plan à
$\ang{45} - \phi/2 = \ang{24.8}$ de $\sigma_1$ vertical, de \omterm{def:b1:geological-maps:attitude}{pendage}
\ang{65.2}.

\textbf{16.} $S = \sigma_v' = \qty{83.39}{MPa}$ et
$\sigma_1 - \sigma_3 = 2\times0.85\times83.39/(1.3124 + 0.85) =
\qty{65.55}{MPa}$~; la pente donne $13.11\times5 = \qty{65.55}{MPa}$.

\textbf{17.} $\sigma_3 = 132.44 - 65.55 = \qty{66.88}{MPa}$.

\textbf{18.} La faille glisse quand
$r\sin\ang{120} = 0.85\,(c' + r\cos\ang{120})$ avec
$c' = (83.39 + \sigma_3')/2$ et $r = (83.39 - \sigma_3')/2$~; la résolution donne
$\sigma_3' = \qty{17.18}{MPa}$, $\sigma_3 = \qty{66.23}{MPa}$. C'est plus bas
que pour le meilleur plan~: la faille ne se trouve pas au point où le cercle
touche d'abord la droite, donc le cercle doit encore grandir avant d'atteindre
le point propre de la faille.

\textbf{19.} Chaque \omterm{def:b1:stress-strain-rheology:stress}{contrainte normale} effective, $\sigma_1'$ et $\sigma_3'$
comprises, diminue de $\Delta P$, et leur différence ne change pas~; les
contraintes de cisaillement ne dépendent pas du tout de la pression interstitielle.

\textbf{20.} Le meilleur plan glisse quand la distance égale le rayon~:
$0.85\,(54.67 - \Delta P)/1.3124 = 28.72$, donc
$\Delta P = 54.67 - 28.72\times1.3124/0.85 = 54.67 - 44.34 = \qty{10.33}{MPa}$.

\textbf{21.} La colonne d'eau du puits pèse autant qu'avant, donc la
pompe doit ajouter les \qty{10.33}{MPa} entiers en tête de puits, environ 103 fois
la pression atmosphérique~: le poids d'une colonne d'eau
haute de $\num{10.33e6}/(1000\times9.81) = \qty{1053}{m}$.

\textbf{22.} Avec $\mu = 0.6$ la faille ne pourrait supporter que
$0.6\times40.31 = \qty{24.19}{MPa}$, moins que les \qty{24.87}{MPa} qu'elle
porte~: elle glisserait déjà. Comme elle est calme, son frottement est d'au
moins $0.617$, ou elle a une certaine \omterm{def:b1:stress-strain-rheology:coulomb}{cohésion}.

\textbf{23.} Le plan le mieux orienté glisse quand la pression interstitielle atteint
$49.05 + 10.33 = \qty{59.38}{MPa}$, c'est-à-dire $\lambda = 59.38/132.44 = 0.448$,
bien en deçà de la valeur lithostatique 1.

\textbf{24.} $\Delta P = \sigma_n' - \tau/0.85 = 40.309 - 24.870/0.85 = 40.309 -
29.259 = \qty{11.049}{MPa}$, plus que les \qty{10.33}{MPa} qui feraient bouger un
plan de la meilleure orientation. À une décimale, la hausse de pression interstitielle qui
amène la faille à glisser est
\end{solution}
